\documentclass[10pt]{article}

\usepackage[margin=0.75in]{geometry}
\usepackage{amsmath,amsthm,amssymb,bm}
\usepackage{xcolor}
\usepackage{cancel}
\usepackage{graphicx}
\usepackage{changepage}
\usepackage{circuitikz}
\usepackage{pgfplots}
\usepackage{physics}
\usepackage{tikz}
\usepackage{tikz-3dplot}
\usepackage{hyperref}
\usepackage{siunitx}
\usepackage{fontspec}
\usepackage{relsize}
\usepackage{subfig}
\usepackage{todonotes}
\usepackage{multicol, multirow, booktabs}
\usepackage[breakable]{tcolorbox}
\usepackage[inline]{enumitem}

\theoremstyle{definition}
\newtheorem{problem}{Problem}
\newtheorem{soln}{Solution}

\pgfplotsset{compat=newest}
\usetikzlibrary{lindenmayersystems}
\usetikzlibrary{arrows}
\usetikzlibrary{calc}
\usetikzlibrary{positioning, fit}
\usetikzlibrary{3d, perspective}
\usetikzlibrary{math}
\usetikzlibrary{fadings}
\usetikzlibrary{angles,quotes} 

\definecolor{incolor}{HTML}{303F9F}
\definecolor{outcolor}{HTML}{D84315}
\definecolor{cellborder}{HTML}{CFCFCF}
\definecolor{cellbackground}{HTML}{F7F7F7}
\newcommand{\ui}{\hat{i}}
\newcommand{\uj}{\hat{j}}
\newcommand{\uk}{\hat{k}}
\newcommand{\ux}{\hat{x}}
\newcommand{\uy}{\hat{y}}
\newcommand{\uz}{\hat{z}}
\newcommand{\uv}[1]{\hat{\bv{{#1}}}}
\newcommand{\pr}[1]{{#1}^\prime}
\pgfdeclarelayer{background}  
\pgfsetlayers{background,main}
\AtBeginDocument{\RenewCommandCopy\qty\SI}
\newcommand{\justif}[2]{&{#1}&\text{#2}}

\makeatletter
\newcommand{\boxspacing}{\kern\kvtcb@left@rule\kern\kvtcb@boxsep}
\makeatother
\newcommand{\prompt}[4]{
    \ttfamily\llap{{\color{#2}[#3]:\hspace{3pt}#4}}\vspace{-\baselineskip}
}

\newcommand{\thevenin}[2]{
  \begin{center}
    \begin{circuitikz} \draw
      (0,0) -- (2,0) to[battery1, l_=$V_{Th}\eq#1$] (2,2) 
      to[resistor, l_=$R_{Th}\eq#2$] (0,2)
      ;
      \draw [o-] (-.07,2.079);
      \draw [o-] (-.07,0.079);
    \end{circuitikz}
  \end{center}
}

\newcommand{\norton}[2]{
  \begin{center}
    \begin{circuitikz} \draw
      (0,0) -- (3,0) to[american current source, l_=$I_{N}\eq#1$] (3,2) -- (0,2) (2,0)
      to[resistor, l=$R_{N}\eq#2$] (2,2)
      ;
      \draw [o-] (-.07,2.079);
      \draw [o-] (-.07,0.079);
    \end{circuitikz}
  \end{center}
}

\DeclareMathOperator{\Div}{div}
\DeclareMathOperator{\Curl}{curl}

\DeclareSIUnit\mile{m}

\newlength\stdbls % standard baselineskip
\AtBeginDocument{%
  \setlength\stdbls{\baselineskip}%
  \gdef\stdblsn{\strip@pt\stdbls}%
} 

\newfontface{\Kaufmann}{Kaufmann}
\DeclareTextFontCommand{\kf}{\Kaufmann}
\newcommand{\scriptr}{\fontsize{12pt}{\stdbls}\kf{r}\fontsize{10pt}{\stdbls}}

\newfontface{\KaufmannB}{Kaufmann Bold}
\DeclareTextFontCommand{\kfb}{\KaufmannB}
\newcommand{\bscriptr}{\fontsize{12pt}{\stdbls}\kfb{r}\fontsize{10pt}{\stdbls}}

\newcommand{\highlight}[1]{\colorbox{yellow}{$\displaystyle #1$}}

\newcommand{\ti}[1]{\widetilde{#1}}

\newcommand{\bv}[1]{\bm{\mathrm{#1}}}

\title{Physics 3130H: Assignment V}
\author{Jeremy Favro (0805980) \\ Trent University, Peterborough, ON, Canada}
\date{\today}

\begin{document}
\maketitle

% PROBLEM 1
\begin{problem}
A muon has a mean lifetime of $\qty{2.2}{\micro\second}$ in in its reference frame. Suppose muons are
travelling at $0.92c$ relative to Earth, what is the mean distance a muon will travel as
measured by an observer on Earth?
\end{problem}
\begin{soln}
  Since the muons experience $\Delta t'=\qty{2.2}{\micro\second}$ of subjective time, the Earth observer will experience a time $\Delta t$ given by
  $$\Delta t = \Delta t'\frac{1}{\sqrt{1-0.92^2}}\approx\qty{5.6}{\micro\second}.$$
  At $0.92c$ this means the muons will have travelled a distance of around $\qty{1.5}{\kilo\meter}$.
\end{soln}

% PROBLEM 2
\begin{problem}
An experimenter determines that a particle created at one end of the laboratory apparatus
moved at $0.94c$ and survived for $\qty{0.032}{\micro\second}$, decaying just as it reached the other end.
\begin{enumerate}[label=(\alph*)]
  \item According to the experimenter, how far did the particle move?
  \item In its own reference frame, how long did the particle survive?
  \item According to the particle, what was the length of the laboratory apparatus?
\end{enumerate}
\end{problem}
\begin{soln}~
  \begin{enumerate}[label=(\alph*)]
    \item The particle moved $\qty{0.032}{\micro\second}\cdot c\approx \qty{9.6}{\meter}$.
    \item $\Delta t'=\frac{1}{\gamma}\Delta t\approx \qty{11}{\nano\second}$.
    \item The particle lived for $\qty{11}{\nano\second}$ so experienced a subjective length of $\qty{11}{\nano\second}\cdot c \approx \qty{3}{\meter}$.
          This can also be seen from length contraction where the lab frame experiences $\Delta L \approx \qty{9.6}{\meter}$ and the particle experiences length given by
          $\Delta L'=\frac{1}{\gamma}\Delta L\approx \qty{3}{\meter}$.
  \end{enumerate}
\end{soln}

% PROBLEM 3
\begin{problem}
Consider two events that take place at different points in the $k$ system at the same instant $t$. If the
two points are separated by the distance $\Delta x$, show that the two events are not simultaneous in $k'$
system and separated by a time interval $\Delta t'=-v\gamma\Delta x/c^2$.
\end{problem}
\begin{soln}
  We know that the Lorentz transform tells us that
  $$x'=\gamma(x-vt)$$
  and that
  $$x'=ct'.$$
  Combining these two we see that
  $$ct'=\gamma(x-vt)$$
  so
  $$t'=\gamma(\frac{x}{c}-\frac{vt}{c})$$
  and again from the Lorentz transform we know that $x=ct$ so
  $$t'=\gamma(t-\frac{vct}{c^2})=\gamma(t-\frac{vx}{c^2}).$$
  Since the two events occur in the $k$ frame at the same $t$,
  $$\Delta t'=\gamma(t-\frac{vx_{1}}{c^2})-\gamma(t-\frac{vx_{2}}{c^2})=\gamma (-v\Delta x/c^2)$$
  where $x_1,2$ are the positions of the events in the $k$ frame.
\end{soln}

% PROBLEM 4
\begin{problem}
A star is known to be moving away from Earth at a speed of $\qty{4e4}{\meter\per\second}$. This speed is measured
by measuring the shift of H$\alpha$ ($\lambda=\qty{656.3}{\nano\meter}$). By how much and what direction is the shift of H$\alpha$
line?
\end{problem}
\begin{soln}
  This is a source receding relativistic doppler problem, hence we use
  $$\lambda_{\mathrm{doppler}}=\lambda_{\mathrm{H}\alpha}\sqrt{\frac{1+v/c}{1-v/c}}\approx \qty{656.4}{\nano\meter}.$$
  The effect is very subtle since $\qty{4e4}{\meter\per\second}\ll c$ but this is indeed the expected redshift of an object moving away
  from an observer.
\end{soln}

% PROBLEM 5
\begin{problem}
Show that relativistic form of Newton's second law becomes
$$F=m\left(1-\frac{u^2}{c^2}\right)^{-3/2}\frac{du}{dt}$$
\end{problem}
\begin{soln}
  Recall in class we ``redefined'' momentum to be
  $$p=\gamma mu.$$
  Differentiating this to obtain the force,
  $$F=\frac{dp}{dt}=\gamma m\frac{du}{dt}+mu\frac{d\gamma}{dt}.$$
  We can use chain rule to expand
  $$\frac{d\gamma}{dt}=\frac{d\gamma}{du}\frac{du}{dt}$$
  and
  $$\frac{d\gamma}{dv}=\frac{u}{c^2}\gamma^3$$
  so
  \begin{align*}
    F & =\gamma m\frac{du}{dt}+\gamma m\gamma^3\frac{u^2}{c^2}\frac{du}{dt} \\
      & =\gamma m(1+\frac{u^2}{c^2}\gamma^2)\frac{du}{dt}                   \\
      & =\gamma m(1+\frac{\frac{u^2}{c^2}}{1-\frac{u^2}{c^2}})\frac{du}{dt} \\
      & =\gamma m(1+\frac{u^2}{c^2-u^2})\frac{du}{dt}                       \\
      & =m(\frac{c^2}{(c^2-u^2)(1-u^2/c^2)^{1/2}})\frac{du}{dt}             \\
      & =m(\frac{1}{(1-u^2/c^2)(1-u^2/c^2)^{1/2}})\frac{du}{dt}             \\
      & =m((1-u^2/c^2)^{-3/2})\frac{du}{dt}.
  \end{align*}
\end{soln}
\newpage

% PROBLEM 6
\begin{problem}
A common unit of energy used in atomic and nuclear physics is the electron volt ($\unit{\electronvolt}$), the energy
acquired by an electron in falling through a potential difference of one volt. In these units
the mass of the electron is $m_ec^2=\qty{0.511}{\mega\electronvolt}$ and that of the
proton is $m_pc^2=\qty{938}{\mega\electronvolt}$. If their momentum is given by $\qty{100}{\mega\electronvolt}/c$, calculate the following
\begin{enumerate}[label=(\roman*)]
  \item Kinetic Energy
  \item Quantities $\beta$ and $\gamma$ for electron and proton. Identify the one that will behave relativistic
        and another one which behave non-relativistic.
\end{enumerate}
\end{problem}
\begin{soln}
  \begin{enumerate}[label=(\roman*)]
    \item We can find kinetic energy using the general form for relativistic energy,
          $$E=\sqrt{(mc^2)^2+(pc)^2}$$
          and then subtracting $mc^2$.

          For the electron,
          $$E_e=\sqrt{(\qty{0.511}{\mega\electronvolt})^2+(\qty{100}{\mega\electronvolt})^2}\approx \qty{1e8}{\electronvolt}$$
          so
          $$E_{ke}\approx \qty{1e8}{\electronvolt}-mc^2\approx\qty{9.9e7}{\electronvolt}.$$
          For the proton
          $$E_p=\sqrt{(\qty{938}{\mega\electronvolt})^2+(\qty{100}{\mega\electronvolt})^2}\approx \qty{9.4e8}{\electronvolt}$$
          so
          $$E_{kp}\approx \qty{9.4e8}{\electronvolt}-mc^2\approx\qty{5.3e6}{\electronvolt}.$$
    \item In class we found that the total energy of a particle is given by
          $$E=\gamma mc^2\implies \gamma = \frac{E}{mc^2}$$
          and $\beta=v/c$ is easily found by recalling $p=\gamma mv=\gamma mc\beta$ so,
          \begin{center}
            \begin{tabular}{c | c c }
                       & $\gamma$ & $\beta$ \\
              \cline{1-3}
              Electron & $196$    & $0.99$  \\
              Proton   & $1$      & $0.1$
            \end{tabular}
          \end{center}
          evidently the proton is nonrelativistic as $\beta$ is not close to 1 and the electron is relativistic since $\beta \approx 1$.
  \end{enumerate}
\end{soln}
\end{document}